\section{Fabric and topology}
\label{sec:fabric}

\begin{question}
  \textbf{How much does the network limit training?} Does placement on leaf switches matter (same leaf
  vs.\ across leaves, nodes per switch, packed vs.\ balanced)? How does Rorqual's single HCA per node
  compare with clusters that have one per GPU?
\end{question}

\paragraph{Raw bandwidth.} \Cref{fig:comm} places each communication scope on one axis, from NVLink inside a
node to all-to-all across every leaf. On Rorqual, every all-reduce that leaves a node is capped near the
line rate of its single 200\,Gb/s HCA (25\,GB/s). The scopes are those the health checks measure:
per node, per pair, and all the nodes of a check job. The window needs two health-check runs to cover
both one leaf and the whole machine: \texttt{--job-per switch} (default) and \texttt{--job-per cluster}.
Intermediate scopes (2 or 3 leaves) come from the allocation-wide all-reduce of \texttt{tt\_healthcheck}.

\begin{figure}[h]
  \centering\fig[0.75\linewidth]{comm_scopes}
  \caption{Peak bus bandwidth by scope, over the tested message sizes (128\,MiB--1\,GiB). Gray:
    preliminary health checks (\texttt{tt\_healthcheck}); orange: uplink tests between ihpc217 and ihpc316
    (pairs alone, then 7 pairs at once; 2026-10-24); blue: window health check (srygaard). All-to-all
    and all-reduce bus bandwidths are not comparable (different normalization). \syntheticnote{comm_scopes}}
  \label{fig:comm}
\end{figure}

\placeholder{per-scope medians; the drop across leaves including ihpc316 (19.2 vs.\ 24.6\,GB/s in the
NCCL tests, $-35\,\%$ for 217+316) and whether it comes from ihpc316's uplinks.}

\paragraph{Placement at scale: balance or pack?} Up to 12 nodes, placement is analysed with scale in
\cref{sec:scaling} (\cref{fig:placement-scaling}). Beyond 12 nodes no switch holds the job, and the
allocation choice is between spreading it 3 per switch (balanced, over the 3-node leaves) and packing it
into the fewest switches (the big leaves first). \Cref{fig:placement-at-scale} plots throughput per GPU
against total throughput, so both read at once: runs with the same GPU count lie on one ray through the
origin. Along a ray, placement changes only throughput per GPU; from one ray to the next, adding nodes
lowers throughput per GPU but raises the total.

\begin{figure}[h]
  \centering\fig{placement_at_scale}
  \caption{Rorqual beyond one switch: throughput per GPU vs.\ total throughput, for 3 per switch
    (balanced) and fewest switches (packed). Gray rays: same GPU count. \syntheticnote{placement_at_scale}}
  \label{fig:placement-at-scale}
\end{figure}

\placeholder{measured so far, all preliminary: H6 = S6 within 0.2\,\% on Rorqual, repeats within 0.1\,\%; but
H6 crossed the spine outside our control while S6 stayed in one reserved switch, so the pair mixes placement
with possible traffic from other users. P12 split 6+6 across ihpc218 and ihpc316 is 4.5\,\% below S6, not
separable from 6$\to$12-node scaling.}

\paragraph{Per-GPU batch.} With FSDP, every step all-gathers the parameters and reduce-scatters the
gradients: a volume that does not depend on the batch. If communication dominates, more tokens per GPU
per step cost almost no extra time (\cref{fig:batch}).

\begin{figure}[h]
  \centering\fig{batch_effect}
  \caption{Throughput vs.\ tokens per GPU per step. On 2 nodes, all runs use BF16 AdamW moments (low-memory
    variant).}
  \label{fig:batch}
\end{figure}

\placeholder{Rorqual 2 nodes: $1\to2$ sequences of 8192 gives $+80\,\%$ (954 $\to$ 1722); Fir 3 nodes:
$+24\,\%$. Step time on Rorqual barely changes with the batch: communication-bound.}
